\section{Socket-Heating Characterization}
\label{sec:ch4-socket-heating-characterization}

The socket-heating characterization compared a normal socket with a degraded
socket under similar current levels. The purpose was to verify the Chapter 3
assumption that socket temperature is not determined by current alone; contact
condition changes the rate and final magnitude of heating.
The plotted trends in this section were generated from the full
characterization records in Appendix~\ref{app:c}.

The water-heater-only characterization trend is shown in
Figure~\ref{fig:ch4-char-water-heater}.
\begin{figure}[H]
  \centering
  \socketCharPlot{1}{240}
  \caption{Water heater temperature trend during socket-heating characterization.}
  \label{fig:ch4-char-water-heater}
\end{figure}

Figure~\ref{fig:ch4-char-water-heater} shows that at about 5.2A, both socket
conditions stayed below the 38.2°C comparison point within the 240min
window. The corroded socket still ended hotter than the normal socket, showing
the beginning of contact-condition influence even at the lowest tested load.

The water heater plus one halogen lamp trend is shown in
Figure~\ref{fig:ch4-char-one-lamp}.
\begin{figure}[H]
  \centering
  \socketCharPlot{2}{240}
  \caption{Water heater plus one halogen lamp temperature trend during characterization.}
  \label{fig:ch4-char-one-lamp}
\end{figure}

Figure~\ref{fig:ch4-char-one-lamp} shows clearer separation between normal and
corroded behavior. The corroded socket reached the 38.2°C comparison
point, while the normal socket did not reach that point within the same window.

The water heater plus two halogen lamps trend is shown in
Figure~\ref{fig:ch4-char-two-lamps}.
\begin{figure}[H]
  \centering
  \socketCharPlot{3}{240}
  \caption{Water heater plus two halogen lamps temperature trend during characterization.}
  \label{fig:ch4-char-two-lamps}
\end{figure}

Figure~\ref{fig:ch4-char-two-lamps} shows that the degraded socket reached the
comparison temperature much earlier than the normal socket. The terminal
temperature also rose above the NTC temperature, supporting the need to treat
the external NTC as an under-reading proxy rather than the actual hotspot.

The water heater plus three halogen lamps trend is shown in
Figure~\ref{fig:ch4-char-three-lamps}.
\begin{figure}[H]
  \centering
  \socketCharPlot{4}{240}
  \caption{Water heater plus three halogen lamps temperature trend during characterization.}
  \label{fig:ch4-char-three-lamps}
\end{figure}

Figure~\ref{fig:ch4-char-three-lamps} shows the strongest time-to-threshold
contrast in the characterization set: the corroded socket reached the
comparison point after only 7min, while the normal socket reached it after
52min.

The water heater plus four halogen lamps trend is shown in
Figure~\ref{fig:ch4-char-four-lamps}.
\begin{figure}[H]
  \centering
  \socketCharPlot{5}{195}
  \caption{Water heater plus four halogen lamps temperature trend during characterization.}
  \label{fig:ch4-char-four-lamps}
\end{figure}

Figure~\ref{fig:ch4-char-four-lamps} shows that the high-current degraded
trial developed rapidly and was stopped after a short observation window. The
terminal temperature reached 45.8°C within 4min even though the final
NTC reading was only 36.9°C.

Table~\ref{tab:ch4-heating-characterization} summarizes the characterization
data used to interpret these trends.
\begin{table}[H]
  \centering
  \caption{Summary of socket-heating characterization data}
  \label{tab:ch4-heating-characterization}
  \small
  \setlength{\tabcolsep}{2pt}
  \begin{tabular*}{\textwidth}{@{\extracolsep{\fill}}p{0.22\textwidth}p{0.105\textwidth}p{0.13\textwidth}p{0.08\textwidth}p{0.11\textwidth}p{0.135\textwidth}p{0.09\textwidth}@{}}
    \toprule
    Load case & Socket & Time to 38.2°C & Mean A & Final NTC & Final terminal & Window (min) \\
    \midrule
    Water heater & Normal & -- & 5.18 & 33.9 & 31.5 & 240 \\
    Water heater & Corroded & -- & 5.16 & 35.6 & 37.4 & 240 \\
    Water heater + 1 lamp & Normal & -- & 7.12 & 37.7 & 37.6 & 240 \\
    Water heater + 1 lamp & Corroded & 95 & 7.10 & 40.4 & 43.2 & 240 \\
    Water heater + 2 lamps & Normal & 105 & 9.07 & 40.2 & 40.8 & 240 \\
    Water heater + 2 lamps & Corroded & 48 & 9.05 & 41.0 & 44.5 & 120 \\
    Water heater + 3 lamps & Normal & 52 & 11.06 & 43.2 & 45.8 & 240 \\
    Water heater + 3 lamps & Corroded & 7 & 11.15 & 38.8 & 53.0 & 8 \\
    Water heater + 4 lamps & Normal & 44 & 13.00 & 43.4 & 46.4 & 195 \\
    Water heater + 4 lamps & Corroded & -- & 13.10 & 36.9 & 45.8 & 4 \\
    \bottomrule
  \end{tabular*}
\end{table}

Table~\ref{tab:ch4-heating-characterization} confirms that degraded contact
changed both heating speed and hotspot severity. At about 9A, the corroded
socket reached 38.2°C in 48min compared with 105min for the normal
socket. At about 11A, the corroded socket reached the same point in 7min
compared with 52min for the normal socket. These results support the need for
the compensated thermal model described in Section~\ref{sec:ch3-socket-heating-models}.

